\documentclass[10pt,a4paper]{amsart}
\usepackage[margin=19mm]{geometry}
\usepackage{amsmath,amssymb,mathtools}
\usepackage[scr=rsfs,cal=euler]{mathalfa}
\usepackage{xfrac}
\usepackage[unicode,hidelinks,bookmarks=false]{hyperref}
\newcommand{\ie}{i.e.\ }
\newcommand{\cf}{cf.\ }
\newcommand{\resp}{respectively\ }
\newcommand{\Subsection}{\subsection}
\providecommand{\nobreakdash}{\nobreak\hbox{-}}
\setlength{\emergencystretch}{2em}
\multlinegap0pt
\usepackage{bm}
\usepackage{bbm}
\usepackage{dsfont}
\usepackage{xspace}
\usepackage{cancel}
\usepackage{mathtools}
\usepackage{amsthm}
\makeatletter
\@ifundefined{thm}{\newtheorem{thm}{Theorem}}{}
\@ifundefined{lem}{\newtheorem{lem}[thm]{Lemma}}{}
\makeatother
\newcommand{\wics}{\operatorname{wics}}
\title[K3 surfaces of any Artin-Mazur height over $\mathbb{F}\_5$ an]{K3 surfaces of any Artin-Mazur height over $\mathbb{F}\_5$ and $\mathbb{F}\_7$ via quasi-$F$-split singularities and GPU acceleration}
\author{Ryan Batubara, Jack J Garzella, Alex Pan}
\date{Source: arXiv:2502.12428v3 (CC BY 4.0)}
\begin{document}
\maketitle

\noindent\emph{Source and licence.} K3 surfaces of any Artin-Mazur height over $\mathbb{F}_5$ and $\mathbb{F}_7$ via quasi-$F$-split singularities and GPU acceleration. Version arXiv:2502.12428v3 (CC BY 4.0), distributed under the Creative Commons Attribution 4.0 International licence, as recorded in the arXiv licence metadata for this exact version and shown on the abstract page https://arxiv.org/abs/2502.12428v3. Full rights notice: licenses/arxiv-2502.12428-CC-BY-4.0.txt.\par\smallskip

\noindent\textit{Editorial scope.} Counted unit: the Thm whose source label is \texttt{None} in the article (source0.tex lines 1672--1677, source label \texttt{None}), delivered together with the article's own complete original proof of it (source0.tex lines 1679--1707, one \texttt{proof} environment of 29 lines). No mathematics is written, repaired or re-derived: statement and proof are the article's own text line for line. Required context retained uncounted: lem:wics:size (lines 1138--1141). Every definition, display and label of those lines is retained unchanged. Omitted from this package and not redistributed: 1--1137, 1142--1671, 1678--1678, 1708--2195. Nothing inside a retained excerpt is omitted or altered: every retained body is delivered exactly as the article's lines are written, so no omission receipt of any kind accompanies this package. Citations: the retained excerpts carry no citation command at all, so no bibliography entry of the article is transcribed and no reference is redistributed. Reference adaptation: none was needed and none was made; every reference inside the delivered unit resolves inside it, so no unresolved reference or citation is produced. Printed slips kept literal: no printed word, symbol, hypothesis or number of the counted statement or of its proof was altered. Layout and class: the article's own class is \texttt{article} with packages including babel, inputenc, fontenc, geometry, mathptmx, amsmath, colonequals, amssymb, gensymb, mathrsfs, amscd, amsthm, mathtools, graphicx; the wrapper is the lane's pinned \texttt{amsart} editorial wrapper at 10pt on A4, which restates the theorem-like environments the retained text needs and adds the article's own macro definitions that the retained text needs (\texttt{\textbackslash wics}), with extra wrapper packages (bm, bbm, dsfont, xspace, cancel, mathtools). The delivered numbering is the unit's own. Assets: no retained span contains a figure environment, a graphics-inclusion command, a PDF-page inclusion, a table or a TikZ picture; no image file is copied into this package, the package declares no asset dependency and no image file is redistributed with it. Licence: version arXiv:2502.12428v3 of this article is distributed under the Creative Commons Attribution 4.0 International licence, as recorded in the arXiv licence metadata for this exact version and shown on the abstract page https://arxiv.org/abs/2502.12428v3. A full rights notice is delivered at licenses/arxiv-2502.12428-CC-BY-4.0.txt. Characters: every retained excerpt is pure ASCII, so the delivered engine, XeLaTeX with its default Latin Modern text font, reports no missing character.
\begin{lem}
    \label{lem:wics:size}
    $|\wics(k, n)| = \binom{k + n - 1}{n - 1}$
\end{lem}

\begin{thm}
    Let $f \in \mathbb{Z}[x_1, \dots, x_n]$ be homogeneous of degree \(h\), 
    and let $m$ be a non-inclusive upper bound on the coefficients. 
    Then, the coefficients of $f ^ k$ are bounded above by 
    $\left((m - 1) \cdot \binom{h + n - 1}{n - 1}\right)^ k$
\end{thm}

\begin{proof}
    We induct on $k$. The maximum coefficient 
    of $f^0$ is $1$, which is bounded above by 
    
    \noindent$\left((m - 1) \cdot \binom{h + n - 1}{n - 1}\right)^ 0 = 1$.
    Let $(d_1, \dots , d_n)$ denote the exponent 
    tuple of a term of $f^k$, let $(d'_1, \dots , d'_n)$ 
    denote the exponent tuple of a term of 
    $f^{k + 1}$, and let $(a_1, \dots , a_n)$ 
    denote the exponent tuple of a term of $f$.

    Consider an arbitrary term of $f^{k + 1}$. 
    To find all terms in the unreduced 
    expansion of $f^k \cdot g$ that contribute 
    to that term, we look at 
    terms of $f^k$ with degree sequences of 
    the form $(d'_1 - a_1, \dots , d'_n - a_n)$. 
    The number of these terms of $f^k$ is bounded 
    above by $|wics(h, n)|$, or $\binom{h + n - 1}{n - 1}$ 
    by Lemma \ref{lem:wics:size}. Using our 
    inductive assumption, the maximum coefficient 
    of $f^k$ is bounded above by 
    $\left((m - 1) \cdot \binom{h + n - 1}{n - 1}\right)^ k$, 
    so multiplying each of these by coefficients 
    of $f$, which have a maximum value of 
    $m - 1$, and adding up $\binom{h + n - 1}{n - 1}$ 
    copies of these gives 
    $\left((m - 1) \cdot \binom{h + n - 1}{n - 1}\right)^{k + 1}$.
\end{proof}

\end{document}
